%!LW recipe=Docker: XeLaTeX
\documentclass{homework}   % 要顯示圖片 caption：\documentclass[showcaption]{homework}

%---------------------------------------------------------------
% 作業資訊（只需改這裡）
%---------------------------------------------------------------
\title{Course Name -- Assignment 1}
\subtitle{你的名字\quad 00000000X}

%---------------------------------------------------------------
% 常用數學巨集（依課程增減）
%---------------------------------------------------------------
\newcommand{\Z}{\mathbb{Z}}
\newcommand{\xor}{\oplus}
\newcommand{\rand}{\xleftarrow{\mathrm{R}}}   % r <-R- Z_p
\DeclareMathOperator{\Adv}{Adv}

\begin{document}
\maketitle

%===============================================================
\problem{題目名稱}
%===============================================================
題目敘述。數學式例如 $a \xor b$、$r \rand \Z_p$。

\begin{enumerate}
  \item 第一小題
  \begin{proof}
    證明內容。
    \[
      \Adv[\mathcal{A}] \le \epsilon. \qedhere
    \]
  \end{proof}

  \item 第二小題
  \begin{align*}
    x^5 + x + 1 &= (x^4 + x^2 + 1)(x) + (x^3 + 1)
  \end{align*}
\end{enumerate}

%===============================================================
\problem{圖、表與程式碼}
%===============================================================
截圖放在 \texttt{figure/}，直接寫檔名即可（例如 \verb|\includegraphics{task1.png}|）。
\begin{figure}[H]
  \centering
  \includegraphics[width=0.5\linewidth]{example-image}
  \caption{圖片說明（預設隱藏）}
\end{figure}

\begin{table}[H]
  \centering
  \begin{tabular}{ccc}
    \toprule
    欄位 A & 欄位 B & 欄位 C \\
    \midrule
    1 & 2 & 3 \\
    \bottomrule
  \end{tabular}
\end{table}

程式碼可以直接寫在文件中，或用 \verb|\lstinputlisting[language=Python]{code/xxx.py}| 引入 \texttt{code/} 裡的檔案。
\begin{lstlisting}[language=Python]
def xor(a: bytes, b: bytes) -> bytes:
    # 逐位元組 XOR
    return bytes(x ^ y for x, y in zip(a, b))
\end{lstlisting}

\end{document}
